% =====================================================================
% Comparación CRM actual vs CRM nuevo
% Compilar:  latexmk -pdf comparacion-crm-nuevo.tex
%
% Todo lo que afirma este documento se verificó EN PANTALLA el
% 2026-09-22, entrando a las dos aplicaciones. Cuando algo no se pudo
% verificar, se dice explícitamente en la sección de límites.
% =====================================================================
\documentclass[11pt,a4paper]{article}

\usepackage[utf8]{inputenc}
\usepackage[T1]{fontenc}
\usepackage[spanish,es-nodecimaldot]{babel}
\usepackage{lmodern}
\usepackage[margin=2.3cm,headheight=14pt]{geometry}
\usepackage{xcolor}
\usepackage{booktabs}
\usepackage{tabularx}
\usepackage{longtable}
\usepackage{array}
\usepackage{enumitem}
\usepackage{titlesec}
\usepackage{fancyhdr}
\usepackage{tcolorbox}
\usepackage[hidelinks]{hyperref}

\definecolor{primary}{HTML}{1F3864}
\definecolor{accent}{HTML}{2E75B6}
\definecolor{soft}{HTML}{EEF3FA}
\definecolor{warn}{HTML}{C0504D}
\definecolor{ok}{HTML}{2E7D32}
\definecolor{gray70}{HTML}{4D4D4D}

\hypersetup{
  pdftitle={Comparacion CRM actual vs CRM nuevo},
  pdfsubject={Inventario funcional y brecha para reemplazar el CRM en produccion},
  pdfkeywords={CRM, whaticket, migracion, canales, comparacion}
}

\pagestyle{fancy}
\fancyhf{}
\fancyhead[L]{\small\color{gray70}Comparación CRM actual vs CRM nuevo}
\fancyhead[R]{\small\color{gray70}\thepage}
\renewcommand{\headrulewidth}{0.4pt}

\titleformat{\section}{\normalfont\Large\bfseries\color{primary}}{\thesection}{0.6em}{}
\titleformat{\subsection}{\normalfont\large\bfseries\color{accent}}{\thesubsection}{0.6em}{}
\setlist[itemize]{leftmargin=1.3em,itemsep=2pt,topsep=3pt}
\setlist[enumerate]{leftmargin=1.5em,itemsep=2pt,topsep=3pt}

\newcolumntype{L}{>{\raggedright\arraybackslash}X}
\newcommand{\cod}[1]{\texttt{\small #1}}
\newtcolorbox{nota}[1][]{colback=soft,colframe=accent,boxrule=0.6pt,
  arc=2pt,left=6pt,right=6pt,top=5pt,bottom=5pt,#1}
\newtcolorbox{alerta}[1][]{colback=white,colframe=warn,boxrule=0.8pt,
  arc=2pt,left=6pt,right=6pt,top=5pt,bottom=5pt,#1}

\newcommand{\si}{\textcolor{ok}{\textbf{Sí}}}
\newcommand{\no}{\textcolor{warn}{\textbf{No}}}
\newcommand{\parcial}{\textcolor{accent}{\textbf{Parcial}}}
\newcommand{\pordef}{\textcolor{gray70}{\textbf{Por definir}}}
\newcommand{\vars}[1]{\cod{\{\{#1\}\}}}

\begin{document}

\begin{titlepage}
\centering
\vspace*{3cm}
{\Huge\bfseries\color{primary} Comparación del CRM\par}
\vspace{0.6cm}
{\Large\color{gray70} Qué tiene el nuevo, qué le falta para reemplazar al actual\par}
\vspace{2cm}
{\large CRM actual: \cod{app.whaticket.com}\par}
\vspace{0.3cm}
{\large CRM nuevo: \cod{staging-crm.zgameslatam.com} --- producto \emph{Ticket} \cod{v0.1.0}\par}
\vspace{0.4cm}
{\large 22 de septiembre de 2026\par}
\vfill
\begin{minipage}{0.82\textwidth}
\small\color{gray70}
Todo lo que afirma este documento se verificó \textbf{entrando a las dos
aplicaciones} el 22 de septiembre de 2026, no leyendo documentación ni código. Cada
afirmación se acompaña de lo que se vio en pantalla.

\vspace{0.5em}
La sección \ref{sec:limites} enumera lo que \textbf{no} se pudo verificar y por qué.
Se prefiere decir ``no lo sé'' antes que dar por buena una suposición: un inventario
con un dato inventado deja de servir para decidir.
\end{minipage}
\vspace{2cm}
\end{titlepage}

\tableofcontents
\newpage

% =====================================================================
\section{El veredicto, en una página}

El CRM nuevo está más avanzado de lo que sugiere su número de versión: tiene
pantallas que el actual no tiene y varias resueltas con mejor criterio. Pero
\textbf{no puede reemplazar al actual hoy}, y el motivo no es una pantalla que falte.

\begin{alerta}
\textbf{El bloqueante es de canales, no de funcionalidad.} El CRM nuevo solo admite
\emph{Web Chat} y \emph{WhatsApp QR}. Facebook, Instagram y WhatsApp Cloud API figuran
como ``Próximamente'' en su propio diálogo de conexión, y \textbf{Telegram no figura
en absoluto}. El CRM actual tiene los cuatro canales conectados y operando.

\vspace{0.4em}
Cruzado con la medición de tráfico ya existente, \textbf{Facebook representa el
32,8\,\% de las conversaciones}. Mientras ese canal no entre, migrar significa perder
un tercio de la operación.
\end{alerta}

\begin{center}
\begin{tabularx}{\textwidth}{@{}lcL@{}}
\toprule
\textbf{Dimensión} & \textbf{Estado} & \textbf{Qué significa} \\
\midrule
Canales de entrada      & \no       & Falta Facebook, Instagram y Telegram \\
Cobertura de pantallas  & \parcial  & Cubre 8 de 15 secciones del actual \\
Chat y mensajería       & \parcial  & Falta el atajo \cod{/} de respuestas rápidas \\
Permisos y roles        & \si       & Ambos lo tienen; el nuevo mejor estructurado \\
Contactos y etiquetas   & \si       & El nuevo es más completo que el actual \\
Departamentos           & \parcial  & Falta horario de atención y filtro por conexión \\
\bottomrule
\end{tabularx}
\end{center}

\begin{nota}
\textbf{Una aclaración que evita una conclusión equivocada.} La regla de negocio
``estos jugadores, dentro de esta clasificación, solo pueden ser atendidos por estos
operadores'' \textbf{no existe en ninguno de los dos sistemas}. No es una deuda del
CRM nuevo: es funcionalidad que habría que construir. Ver la sección
\ref{sec:visibilidad}.
\end{nota}

% =====================================================================
\section{Alcance y método}

Se recorrieron las dos aplicaciones en el navegador, pantalla por pantalla, con estas
restricciones acordadas:

\begin{itemize}
  \item \textbf{CRM actual (producción):} solo lectura. No se envió ningún mensaje, no
        se guardó ningún formulario y \textbf{no se abrió ningún chat}, porque abrirlo
        asigna el ticket al operador que lo abre y eso alteraría la operación real.
  \item \textbf{CRM nuevo (staging):} se interactuó libremente, por tratarse de un
        entorno de prueba.
\end{itemize}

Para el CRM actual se usó la cuenta \cod{Prueba\_datos} con perfil Administrador; para
el nuevo, \cod{willian test} con perfil Administrador.

\begin{nota}
\textbf{El menú prueba que la ruta existe, no que la pantalla funcione.} Durante el
recorrido, la ruta \cod{/departments} del CRM nuevo cargó en blanco en el primer
intento y mostró contenido recién en el segundo. Por eso el inventario de rutas de la
sección \ref{sec:rutas} se acompaña del recorrido pantalla por pantalla de la sección
\ref{sec:funcional}: la primera dice qué existe, la segunda qué funciona.
\end{nota}

% =====================================================================
\section{Inventario de rutas}
\label{sec:rutas}

\subsection{Lo que el CRM nuevo ya cubre}

Ocho de las quince secciones del CRM actual tienen equivalente:

\begin{center}
\begin{tabular}{@{}lll@{}}
\toprule
\textbf{Sección} & \textbf{CRM actual} & \textbf{CRM nuevo} \\
\midrule
Chats                    & \cod{/tickets}        & \cod{/whaticket/chat} \\
Conexiones               & \cod{/connections}    & \cod{/connections} \\
Contactos                & \cod{/contacts}       & \cod{/contacts} \\
Etiquetas                & \cod{/tags}           & \cod{/labels} \\
Respuestas rápidas       & \cod{/fastresponses}  & \cod{/quick-responses} \\
Equipo / Usuarios        & \cod{/users}          & \cod{/users} \\
Departamentos y chatbots & \cod{/queues}         & \cod{/departments} \\
Configuración            & \cod{/settings}       & \cod{/settings} \\
\bottomrule
\end{tabular}
\end{center}

\subsection{Lo que falta}

Siete secciones del CRM actual no tienen equivalente en el nuevo:

\begin{center}
\begin{tabularx}{\textwidth}{@{}llcL@{}}
\toprule
\textbf{Sección} & \textbf{Ruta actual} & \textbf{¿Necesario?} & \textbf{Comentario} \\
\midrule
Wäbot IA            & \cod{/ai}                  & \no      & Asistente de inteligencia artificial \\
Campañas            & \cod{/campaigns}           & \pordef  & Envíos masivos \\
Informes            & \cod{/reports}             & \pordef  & Reportería operativa \\
Mensajes programados& \cod{/scheduled-messages}  & \si      & Envíos diferidos \\
Tokens              & \cod{/tokens}              & \pordef  & Credenciales de API \\
Documentación       & \cod{/docs}                & \no      & Ayuda en producto \\
Soporte             & ---                        & \no      & Canal de soporte del proveedor \\
\bottomrule
\end{tabularx}
\end{center}

\noindent\small\textcolor{gray70}{\textbf{Por definir} significa que todavía falta
información para decidir si la sección debe existir en el CRM nuevo. No es un ``tal
vez'': es una decisión pendiente con responsable, y conviene cerrarla antes de estimar
el esfuerzo de la migración.}
\normalsize

\subsection{Lo que el nuevo tiene y el actual no}

Conviene registrarlo, porque una migración también puede ganar capacidades:

\begin{itemize}
  \item \textbf{Tablero propio} (\cod{/dashboard}).
  \item \textbf{Control de acceso con cinco pantallas dedicadas}: Roles, Permisos,
        Permisos por Rol, Módulos del Sistema y \emph{Constructor de Permisos}.
  \item \textbf{Chats internos} entre miembros del equipo
        (\cod{/internal-chat/conversations}).
  \item \textbf{Mensajes destacados} (\cod{/profile/starred-messages}) y
        \textbf{Tickets bloqueados} (\cod{/locked-tickets}).
  \item Buscador de módulos y alternancia de tema claro/oscuro.
\end{itemize}

% =====================================================================
\section{El bloqueante: canales de entrada}
\label{sec:canales}

Esta es la diferencia que decide si la migración es posible.

\subsection{Lo que admite el CRM nuevo}

Su propio diálogo \emph{Agregar conexión} declara el estado de cada canal:

\begin{center}
\begin{tabularx}{\textwidth}{@{}lcL@{}}
\toprule
\textbf{Canal} & \textbf{Estado} & \textbf{Descripción en pantalla} \\
\midrule
Web Chat            & \si            & Widget de chat para el sitio \\
WhatsApp QR         & \si            & Conexión vía código QR \\
WhatsApp Cloud API  & Próximamente   & API oficial de Meta \\
Facebook            & Próximamente   & Mensajes de la página de Facebook \\
Instagram           & Próximamente   & Mensajes directos de cuenta profesional \\
Telegram            & \no            & \textbf{No figura ni como pendiente} \\
\bottomrule
\end{tabularx}
\end{center}

En staging hay dos conexiones configuradas: una de WhatsApp, \emph{desconectada}, y
una de Web Chat, conectada.

\subsection{Lo que tiene el CRM actual}

Diez conexiones, todas conectadas, sobre cuatro canales distintos:

\begin{center}
\begin{tabular}{@{}lll@{}}
\toprule
\textbf{Conexión} & \textbf{Identificador} & \textbf{Canal} \\
\midrule
SortiOficial                  & 8794286558              & Telegram \\
Sorti365                      & 681388465066532         & Facebook \\
Sorti365 Oficial              & 17841463 / 78834867     & Instagram \\
Prueba ZTicket                & \cod{77d0c521-370e-\ldots} & Web chat / API \\
Plataforma anterior           & 593994038869            & WhatsApp \\
Agentes nueva                 & 593999303548            & WhatsApp \\
JUGADORES PLATAFORMA          & 593965195842            & WhatsApp \\
AGENTES                       & 593984701184            & WhatsApp \\
AGENTES VIP                   & 593988535399            & WhatsApp \\
ALTERNA PLATAFORMA            & 593983902175            & WhatsApp \\
\bottomrule
\end{tabular}
\end{center}

\begin{alerta}
\textbf{Por qué esto manda sobre todo lo demás.} Aunque el CRM nuevo alcanzara la
paridad en las quince pantallas, seguiría sin poder recibir un solo mensaje de
Facebook, Instagram o Telegram. Con Facebook en el 32,8\,\% del tráfico medido, la
migración no es viable hasta que ese canal exista. Telegram, además, ni siquiera está
en la hoja de ruta declarada del producto.
\end{alerta}

% =====================================================================
\section{Comparación funcional, pantalla por pantalla}
\label{sec:funcional}

\subsection{Respuestas rápidas}

El CRM actual describe el mecanismo en su propia pantalla, textualmente:

\begin{nota}
``Mensajes listos que el equipo inserta en el chat escribiendo \textbf{``/'' seguido
del atajo}.''
\end{nota}

Y los atajos en uso lo confirman: \cod{///*}, \cod{/++}, \cod{//000},
\cod{/1 CAMBIO DE CLAVE}, \cod{/20MINUTOS}, \cod{/\$300}, \cod{/3ERA EDAD},
\cod{/A1.1JG}. Son decenas, en uso diario.

\textbf{En el CRM nuevo el atajo no existe.} Se escribió \cod{/} y luego \cod{/sal} en
el campo de mensaje y no se desplegó ningún menú. Las respuestas rápidas se invocan
desde un icono de rayo en la barra de escritura, y al abrirlo con el chat de prueba
respondió: ``No hay respuestas rápidas disponibles para el departamento de este chat''.

\begin{center}
\begin{tabularx}{\textwidth}{@{}Lcc@{}}
\toprule
\textbf{Capacidad} & \textbf{Actual} & \textbf{Nuevo} \\
\midrule
Invocar con \cod{/} desde el chat        & \si & \no \\
Invocar desde un botón del compositor    & --- & \si \\
Buscar por atajo                         & \si & \si \\
Buscar por contenido del mensaje         & \si & \no \\
Filtrar por departamento                 & \si & \no \\
Filtrar por activo / inactivo            & --- & \si \\
Asignar a varios departamentos           & \si & --- \\
Variables de plantilla                   & \si & \si \\
Vista de lista y de tarjetas             & \si & \no \\
\bottomrule
\end{tabularx}
\end{center}

\begin{nota}
\textbf{Dos precisiones para no exagerar la brecha.}

\emph{Primera:} el alcance por departamento \textbf{no es nuevo}. El CRM actual también
tiene columna de departamentos en sus respuestas rápidas, y además admite varios por
respuesta.

\emph{Segunda:} las variables de plantilla \textbf{existen en ambos}, y en el nuevo con
mejor experiencia --- se insertan con fichas: \emph{Saludos}, \emph{Apodo de contacto},
\emph{Nombre de contacto}, \emph{Número de contacto}, \emph{Número de protocolo},
\emph{Correo electrónico}. Lo que cambia son los nombres internos: el actual usa
\vars{greeting}, \vars{contactName} y \vars{contactTreatment}; el nuevo, \vars{greetings}
y \vars{username}. \textbf{Es un mapeo de migración, no una reescritura.}
\end{nota}

Queda un problema práctico: con atajos crípticos como \cod{///*} o \cod{//000}, poder
buscar \emph{solo} por atajo obliga a saberse el código de memoria. En el CRM actual se
puede buscar por el texto del mensaje.

\subsection{Chat: acciones sobre el mensaje y adjuntos}

El menú contextual de un mensaje en el CRM nuevo ofrece: \textbf{Reenviar},
\textbf{Responder}, \textbf{Seleccionar}, \textbf{Destacar} y \textbf{Copiar mensaje}.
La opción \emph{Destacar} conecta con la pantalla de Mensajes destacados.

El botón de adjuntos ofrece \textbf{Imagen}, \textbf{Video} y \textbf{Documento}.

\begin{center}
\begin{tabularx}{\textwidth}{@{}Lcc@{}}
\toprule
\textbf{Capacidad} & \textbf{Actual} & \textbf{Nuevo} \\
\midrule
Copiar mensaje                 & \si & \si \\
Reenviar, responder, destacar  & \si & \si \\
Selección múltiple de mensajes & \si & \si \\
Adjuntar imagen, video, documento & \si & \si \\
Arrastrar y soltar imágenes    & \si & \textit{sin verificar} \\
\bottomrule
\end{tabularx}
\end{center}

\subsection{Departamentos y chatbots}

El CRM actual tiene \textbf{nueve departamentos} en uso (Agente, AGENTES OPERATIVOS,
AGENTES VIP, Departamento de Marketing, JEFE ATC, Jugadores, Jugadores VIP y dos de
prueba), con buscador y \textbf{filtro por conexión}.

\begin{nota}
\textbf{El departamento del CRM actual es una unidad de configuración completa, no una
etiqueta.} Su editor tiene \textbf{seis pestañas}, y el encabezado del modal declara el
modelo: \emph{``Heredando todas las configuraciones globales''}. \textbf{Cada pestaña
tiene su propio interruptor \emph{Sobrescribir configuración global}}, de modo que un
departamento puede apartarse de la configuración general en cualquiera de estos ejes por
separado.
\end{nota}

\begin{center}
\begin{tabularx}{\textwidth}{@{}lL@{}}
\toprule
\textbf{Pestaña} & \textbf{Qué contiene} \\
\midrule
Chatbot & Nombre, color, \textbf{conexiones vinculadas} (el departamento se ata a canales
          concretos), mensaje de saludo con editor enriquecido, pasos del chatbot y
          \textbf{vista previa en vivo} del mensaje sobre la conexión elegida \\
Equipo y asignación & Asignación automática, asignar a usuarios fuera de línea,
          redistribuir chats de usuarios no disponibles, delegar al Wäbot tras la
          asignación; y por agente, \emph{miembro} y \emph{asignación automática} \\
Horario de atención & Activar horario, enviar al final del chatbot, responder con Wäbot
          fuera de horario \\
Resolución del chat & Motivo del contacto al cerrar, resolución automática, mensaje de
          despedida de la conexión, resolver chats inactivos del Wäbot \\
Encuesta de satisfacción & La encuesta completa, configurable \textbf{por departamento} \\
Wäbot IA & Configuración del asistente para ese departamento \\
\bottomrule
\end{tabularx}
\end{center}

El CRM nuevo tiene un departamento cargado y un editor con \textbf{tres} pestañas:
\emph{Departamento} (nombre, color, mensaje de bienvenida con fichas de variables),
\emph{Usuarios y asignación} y \emph{Resolución del chat}.

\begin{center}
\begin{tabularx}{\textwidth}{@{}Lcc@{}}
\toprule
\textbf{Capacidad} & \textbf{Actual} & \textbf{Nuevo} \\
\midrule
Buscar departamento                          & \si & \si \\
Color y mensaje de bienvenida                & \si & \si \\
Asignación automática de chats               & \si & \si \\
Asignar a usuarios fuera de línea            & \si & \si \\
Redistribuir chats de no disponibles         & \si & \si \\
Resolución automática por tiempo             & \si & \si \\
Motivo del contacto al cerrar                & \si & \si \\
Filtrar la lista por conexión                & \si & \no \\
\textbf{Sobrescribir la configuración global} & \si & \parcial \\
\textbf{Conexiones vinculadas al departamento} & \si & \no \\
\textbf{Horario de atención por departamento} & \si & \no \\
\textbf{Encuesta de satisfacción por departamento} & \si & \no \\
\textbf{Pasos de chatbot}                    & \si & \no \\
\textbf{Asistente de IA por departamento}    & \si & \no \\
\bottomrule
\end{tabularx}
\end{center}

\begin{nota}
Los tres interruptores de asignación del CRM nuevo son \textbf{idénticos palabra por
palabra} a los del actual. No son una capacidad nueva: son la misma, reimplementada. El
CRM nuevo sí tiene \emph{Sobrescribir configuración global}, pero solo en la pestaña de
resolución del chat; en el actual está en las cinco pestañas configurables.
\end{nota}

\subsection{Contactos}

Aquí el CRM nuevo es claramente más completo.

\begin{center}
\begin{tabularx}{\textwidth}{@{}Lcc@{}}
\toprule
\textbf{Capacidad} & \textbf{Actual} & \textbf{Nuevo} \\
\midrule
Nombre y apellido desglosados (primero y segundo) & \no & \si \\
Bloqueo de contacto                               & \no & \si \\
Estado activo / inactivo                          & \no & \si \\
Filtro por etiquetas                              & \si & \si \\
Filtro por correo electrónico                     & --- & \si \\
Fecha del último mensaje recibido                 & \si & \no \\
\bottomrule
\end{tabularx}
\end{center}

\subsection{Usuarios y permisos}

\textbf{Ambos sistemas tienen control de acceso basado en roles.} El CRM actual lo
resuelve en una pestaña \emph{Perfiles y permisos} dentro de Equipo, con cuatro
perfiles en producción:

\begin{center}
\begin{tabular}{@{}lrrl@{}}
\toprule
\textbf{Perfil} & \textbf{Permisos} & \textbf{Usuarios} & \textbf{Descripción} \\
\midrule
Admin      & 93 & 8  & Acceso total a la plataforma \\
Usuario    & 28 & 14 & Acceso estándar de atención \\
Supervisor & 36 & 5  & --- \\
MKT        & 10 & 4  & --- \\
\bottomrule
\end{tabular}
\end{center}

Los permisos se agrupan por módulo y acción: Chats (16), Llamadas (7), Conexiones (4),
Contactos (8), Etiquetas (4), Departamentos y chatbots (8), Campañas (4), Informes (2),
Configuración (14) y Tokens de API (3). Dentro de Chats aparecen acciones finas como
\emph{Resolver sin enviar encuesta (CSAT)}, \emph{Pausar cierre automático} y
\emph{Ver todos}.

El CRM nuevo lo resuelve con cinco pantallas dedicadas, incluido un \emph{Constructor
de Permisos}, y distingue dos tipos de usuario (Administrador y Operador) en la pantalla
de Usuarios.

\begin{nota}
\textbf{El control de acceso no es un diferencial del CRM nuevo.} Su estructura es más
explícita, pero el sistema actual ya tiene 93 permisos configurados y en uso. Lo que sí
falta en el nuevo, y el actual tiene, es el registro de \textbf{última vez visto} por
usuario y la \textbf{autenticación en dos factores} por usuario.
\end{nota}

\subsection{Etiquetas}

\begin{center}
\begin{tabularx}{\textwidth}{@{}Lcc@{}}
\toprule
\textbf{Capacidad} & \textbf{Actual} & \textbf{Nuevo} \\
\midrule
Nombre y color                & \si & \si \\
Descripción de la etiqueta    & \no & \si \\
Estado activo / inactivo      & \no & \si \\
Volumen en uso                & decenas & 4 \\
\bottomrule
\end{tabularx}
\end{center}

El CRM actual tiene decenas de etiquetas operativas reales --- ACTIVO, ACUMULA SALDO,
Adulto Mayor, AFILIADO, AGENCIA BLOQUEADA, AGENCIA CERRADA POR EL AGENTE, entre muchas
otras --- que representan vocabulario de negocio acumulado y deberán migrarse.

% =====================================================================
\section{La regla de visibilidad que no existe en ninguno}
\label{sec:visibilidad}

Se buscó específicamente la capacidad de \textbf{clasificar jugadores en un grupo y
restringir qué operadores pueden verlos o atenderlos}. Se revisaron, en ambos sistemas,
las pantallas de Departamentos, Contactos, Etiquetas y Usuarios, y el detalle de un
perfil de permisos del CRM actual.

\begin{alerta}
\textbf{Esa regla no existe en ninguno de los dos sistemas.} Los mecanismos más
cercanos son:

\begin{itemize}
  \item \textbf{CRM actual:} el permiso \emph{Ver todos} dentro de Chats, que es binario
        --- el operador ve todos los chats o solo los suyos --- y la pertenencia a un
        departamento. No hay forma de decir ``este grupo de contactos solo lo ve este
        grupo de operadores''.
  \item \textbf{CRM nuevo:} la asignación de usuarios a un departamento con reparto
        automático de chats. Agrupa \emph{operadores}, no \emph{contactos}.
\end{itemize}

Las etiquetas de contacto, en los dos sistemas, \textbf{clasifican pero no restringen}.
\end{alerta}

Conclusión: si el negocio necesita esta regla, es \textbf{desarrollo nuevo}, y conviene
plantearla como requisito del CRM nuevo mientras está en construcción, en lugar de
esperar que aparezca sola al alcanzar la paridad.

% =====================================================================
\section{Estados de conversación y asignación}
\label{sec:estados}

ATC trabaja con los estados de la conversación y con el filtro de responsable: cada
operador ve por defecto solo lo suyo, para que el trabajo de los demás no se le
amontone en la lista. Esta sección documenta ese mecanismo, qué guarda el ETL de él y
qué tiene el CRM nuevo.

\subsection{Los cinco estados del CRM actual}

La lista de chats ofrece cinco filtros. Dos están ocultos tras un botón \cod{+2}:

\begin{center}
\begin{tabularx}{\textwidth}{@{}llrL@{}}
\toprule
\textbf{Filtro} & \textbf{Columna del ETL} & \textbf{Filas} & \textbf{Observación} \\
\midrule
En atención & \cod{tickets.status='open'}    & 13     & \\
En espera   & \cod{tickets.status='pending'} & 1      & \\
Resueltos   & \cod{tickets.status='closed'}  & 25.539 & \\
Aplazados   & \cod{tickets.snoozed\_until}   & \textbf{0} & \textbf{Nunca se capturó uno} \\
Todos       & ---                            & ---    & No es un estado \\
\bottomrule
\end{tabularx}
\end{center}

\begin{alerta}
\textbf{El 51,2\,\% de \cod{tickets.status} está en NULL, y no es un estado: es un hueco
de captura.} Son 26.818 filas, y \textbf{todas cortan el 30 de junio de 2026} --- la
fecha en que arrancó el monitor. Los tickets anteriores nunca tuvieron su detalle
\cod{/tickets/:id} traído. El corte es idéntico en las dos cuentas: \cod{sistemas} con
24.279 filas y \cod{datos} con 2.539.
\end{alerta}

\subsection{El filtro de responsable}

El desplegable de la cabecera tiene dos niveles:

\begin{enumerate}
  \item \textbf{Míos} --- seleccionado por defecto --- o \textbf{Equipo}.
  \item Dentro de Equipo, una sección \textbf{Responsable} con \emph{Sin usuario} y cada
        operador por nombre.
\end{enumerate}

Esto se corresponde con el permiso \textbf{Ver todos} dentro del módulo Chats: el perfil
``Usuario'' tiene 11 de los 16 permisos de Chats, y quien no tenga ese permiso no puede
salir de ``Míos''. La hipótesis de que la visibilidad se condiciona por rol es correcta.

\subsection{El CRM nuevo no tiene estados}

Verificado sobre el árbol de accesibilidad, no solo sobre una captura. Su lista de chats
ofrece: filtrar por agente, buscar por usuario o nombre, mostrar etiquetas, y un panel
\emph{Filtros} que contiene únicamente \textbf{Etiquetas} y \textbf{Periodo}.

\textbf{No existe ningún filtro por estado de la conversación.} Es el mecanismo central
con el que ATC organiza su día y hoy no tiene equivalente.

\subsection{El historial de asignación se está perdiendo}

\begin{alerta}
El payload de \cod{/tickets/:id} incluye, dentro de \cod{currentConversation}, un array
\textbf{\cod{participations}} con un objeto por tramo de atención: \cod{userId},
\cod{startedAt}, \cod{endedAt} y \cod{firstSentMessageAt}. \textbf{Es el historial
completo de quién tuvo la conversación y entre qué momentos.}

\vspace{0.4em}
Una búsqueda de \cod{participations} en todo el código del ETL --- \cod{monitor/},
\cod{src/} y \cod{db/schema.sql} --- devuelve \textbf{cero resultados}. No está entre las
26 claves que lee el proyector, no hay columna y no hay tabla. Se descarta entero.

\vspace{0.4em}
Y no es solo que se ignore: \textbf{\cod{currentConversation} solo viaja mientras el
ticket está abierto}. Medido, hay 13 conversaciones con historial en \cod{raw\_tickets},
que son exactamente los 13 tickets en estado \cod{open}. Al cerrarse el ticket el array
desaparece del payload. \textbf{Lo que no se capture en vivo se destruye.}
\end{alerta}

Sobre esas 13 conversaciones: 18 tramos de asignación en total, un máximo de 4 operadores
en una sola conversación, promedio de 1,38, y \textbf{3 de 13 (23,1\,\%) pasaron por más
de un operador}.

\begin{nota}
\textbf{Por qué esto importa más allá del inventario.} Ese 23,1\,\% es el mismo fenómeno
que ya se midió por otro camino: \textbf{25.338 calificaciones (24,7\,\% del total)
pertenecen a sesiones con más de un episodio}, donde el departamento, el canal y el
operador se heredan de la conversación de \emph{entrada} en lugar de la que realmente
atendió la interacción calificada. Los rankings por operador atribuyen a la persona
equivocada en una de cada cuatro filas.

\vspace{0.4em}
\cod{participations} contiene exactamente el dato que falta para corregirlo, y está
pasando por el ETL cada 20 segundos.
\end{nota}

% =====================================================================
\section{La encuesta de satisfacción}
\label{sec:csat}

\subsection{Existe en los dos, y el nuevo la replica con los mismos campos}

Conviene decirlo de entrada porque cambia la conclusión: \textbf{el CRM nuevo también
tiene encuesta de satisfacción}, en Configuración \textrightarrow{} Atención. No es una
funcionalidad exclusiva del actual.

\begin{center}
\begin{tabularx}{\textwidth}{@{}LLc@{}}
\toprule
\textbf{Campo en el CRM nuevo} & \textbf{Equivalente en el actual} & \textbf{Estado} \\
\midrule
Activar encuesta CSAT                    & \cod{enabled}            & \si \\
Enviar también en resolución automática  & \cod{sendOnAutoResolve}  & \si \\
Enviar como lista interactiva            & \cod{useInteractiveList} & Próximamente \\
Solicitar comentario adicional           & \cod{requestComment}     & \si \\
Mensaje de solicitud                     & \cod{requestMessage}     & \si \\
Mensaje de agradecimiento                & \cod{thankYouMessage}    & \si \\
Mensaje personalizado por calificación   & \cod{commentBuckets}     & \si \\
\midrule
\textbf{Configuración por departamento}  & \textbf{Sí en el actual} & \no \\
\bottomrule
\end{tabularx}
\end{center}

\begin{alerta}
\textbf{La única diferencia de fondo está en la última fila, y es importante.} En el CRM
actual la encuesta se puede configurar \textbf{por departamento}, con su propio
interruptor de \emph{sobrescribir configuración global}. En el CRM nuevo la encuesta vive
únicamente en la configuración general: \textbf{o está encendida para todos, o para
nadie}.

\vspace{0.4em}
Si se enciende tal como está, el CRM nuevo no podrá evitar encuestar a los agentes ---
que hoy son el 81,7\,\% de los envíos.
\end{alerta}

\begin{nota}
\textbf{Dos hechos que convierten esto en una oportunidad.} Primero, en el CRM nuevo la
encuesta está \textbf{apagada}: todavía no se envió ninguna. Segundo, \emph{enviar como
lista interactiva} figura como \textbf{Próximamente}, así que hoy el nuevo la enviaría
como \textbf{texto plano} --- que es peor que el actual y agrava justamente el problema
de reapertura documentado más abajo.

\vspace{0.4em}
Es decir: los defectos medidos del CRM actual se pueden corregir \textbf{antes} de
encender la encuesta en el nuevo, en lugar de heredarlos.
\end{nota}

\subsection{Cómo funciona hoy}

La encuesta se dispara \textbf{al cerrar el chat}, en milisegundos, y viaja
\textbf{como un mensaje más} por el canal del cliente --- WhatsApp, Telegram, Facebook o
Instagram. No es un formulario aparte: llega a la misma conversación de la persona, en
su propia aplicación de mensajería.

Si el cliente puntúa 5, la encuesta se cierra sola. Si puntúa de 1 a 4, el sistema le
pide un comentario y queda esperando.

\subsection{Problemas registrados, con la medición}

\subsubsection*{1. Es repetitiva}

\begin{center}
\begin{tabular}{@{}lr@{}}
\toprule
Contactos encuestados          & 919 \\
Encuestas enviadas             & 5.164 \\
\textbf{Encuestas por contacto}& \textbf{5,62} \\
\textbf{Máximo a una persona}  & \textbf{167} \\
\bottomrule
\end{tabular}
\end{center}

No hay período de espera entre encuestas: cada cierre dispara una nueva. Una sola persona
recibió 167. En una medición anterior el promedio era 2,80 y el máximo 52, de modo que
\textbf{el problema se duplicó}.

\subsubsection*{2. Se puede focalizar, pero nadie la focalizó}

\begin{center}
\begin{tabularx}{\textwidth}{@{}Lrrr@{}}
\toprule
\textbf{Departamento} & \textbf{Enviadas} & \textbf{Respondidas} & \textbf{Tasa} \\
\midrule
Agente         & \textbf{4.219} & 70 & 1,7\,\% \\
Jugadores      & 679            & 22 & 3,2\,\% \\
AGENTES VIP    & 160            & \textbf{0} & 0,0\,\% \\
Jugadores VIP  & 92             & 1  & 1,1\,\% \\
Sin departamento & 14           & 0  & 0,0\,\% \\
\bottomrule
\end{tabularx}
\end{center}

\textbf{El 81,7\,\% de las encuestas se envía al departamento de agentes}, no a
jugadores. AGENTES VIP recibió 160 y no respondió ninguna.

\begin{alerta}
\textbf{Esto no es una limitación del producto: es una configuración que nadie hizo.} El
editor de departamento del CRM actual tiene una pestaña \emph{Encuesta de satisfacción}
con un interruptor \textbf{Sobrescribir configuración global} --- ``esta configuración se
aplicará solo a este departamento''. La capacidad de dirigir la encuesta existe y está
disponible hoy.

\vspace{0.4em}
Se inspeccionó el departamento \emph{Jugadores} y \textbf{todos sus interruptores de
encuesta están apagados}, de modo que hereda la configuración global. Lo mismo ocurre en
el resto: por eso la encuesta llega por igual a agentes, jugadores VIP y jugadores
nuevos.

\vspace{0.4em}
\textbf{La corrección es de configuración, no de desarrollo}: activar \emph{sobrescribir
configuración global} en los departamentos de agentes y apagar allí la encuesta. Es el
arreglo más barato de todo este documento.
\end{alerta}

\subsubsection*{3. Responder la encuesta reabre la conversación}

\begin{center}
\begin{tabularx}{\textwidth}{@{}Lrrr@{}}
\toprule
& \textbf{Conversaciones} & \textbf{Con actividad posterior} & \textbf{Porcentaje} \\
\midrule
No respondió la encuesta & 5.071 & 243 & 4,8\,\% \\
\textbf{Sí respondió}    & 93    & 24  & \textbf{25,8\,\%} \\
\bottomrule
\end{tabularx}
\end{center}

Se midió la actividad posterior como mensajes llegados más de diez minutos después de la
resolución, para descartar el intercambio propio de la encuesta.

\begin{alerta}
\textbf{Responder la encuesta multiplica por 5,4 la probabilidad de que la conversación
siga viva.} Es consecuencia directa de enviarla como mensaje dentro del canal del
cliente: la respuesta llega como un mensaje entrante y el ticket vuelve a la cola.

\vspace{0.4em}
Una revisión anterior había concluido que la encuesta no reabre conversaciones. Esa
conclusión era incorrecta y queda corregida aquí con la medición.
\end{alerta}

\subsubsection*{4. El promedio que se reporta está sesgado por construcción}

\begin{center}
\begin{tabular}{@{}lrrr@{}}
\toprule
\textbf{Nota} & \textbf{Notas} & \textbf{Cerradas} & \textbf{Trabadas esperando comentario} \\
\midrule
5 & 411 & \textbf{411 (100\,\%)} & 0 \\
4 & 50  & 5  & 45 \\
3 & 8   & 1  & 7 \\
2 & 8   & 3  & 5 \\
1 & 21  & 2  & \textbf{19} \\
\bottomrule
\end{tabular}
\end{center}

Quien puntúa 5 cierra la encuesta de inmediato; quien puntúa bajo recibe una solicitud de
comentario y, si no comenta, queda trabado para siempre. \textbf{68 de las 87 notas malas
(78\,\%) quedan fuera del grupo de encuestas terminadas.}

Promediar únicamente las terminadas da \textbf{4,94}. El promedio real sobre las 498
notas es \textbf{4,65}.

\subsubsection*{5. La tasa de respuesta se calcula mal con facilidad}

La tasa real es del \textbf{9,1\,\%} (498 notas sobre 5.465 encuestas enviadas).
Calcularla dividiendo por \cod{csat\_sent\_at} da entre 1,3\,\% y 2,5\,\%, porque esa
columna estaba en NULL en el 97,4\,\% de las encuestas terminadas. El defecto de captura
ya se corrigió, pero \textbf{solo hacia adelante}: lo perdido no se recupera.

\subsubsection*{6. El comentario contiene datos personales}

El campo de comentario libre es la única fuente de la voz literal del cliente, y trae
números de cédula. Entra en el mismo alcance que el cuerpo de los mensajes para cualquier
política futura de retención o enmascarado.

\subsection{Qué habría que mejorar antes de encenderla en el CRM nuevo}

\begin{enumerate}
  \item \textbf{Configurarla por departamento, que ya se puede.} En el CRM actual basta
        con activar \emph{sobrescribir configuración global} y apagar la encuesta en los
        departamentos de agentes. En el CRM nuevo la encuesta es \textbf{solo global}: si
        se quiere esa focalización, hay que pedirla como requisito antes de encenderla.
  \item \textbf{Un período de espera por contacto.} Sin él, 5,62 encuestas por persona y
        un máximo de 167.
  \item \textbf{Que responder no reabra la conversación.} Es el defecto más caro: afecta
        la cola de trabajo real de ATC, no solo la métrica.
  \item \textbf{Que la nota se registre aunque no haya comentario.} Mientras el
        comentario sea obligatorio para cerrar, el 78\,\% de las notas malas seguirá sin
        contabilizarse.
  \item \textbf{Buscar una integración fuera del hilo de mensajería.} Enviarla como
        mensaje dentro de WhatsApp, Telegram, Facebook o Instagram es la causa común de
        la reapertura y de la sensación de acoso. Un enlace a un formulario externo, o la
        lista interactiva nativa del canal, evitan ambas cosas.
\end{enumerate}

% =====================================================================
\section{Qué falta para reemplazar al CRM actual}

Ordenado por lo que bloquea, no por lo que cuesta.

\begin{enumerate}
  \item \textbf{Canales de entrada.} Facebook, Instagram y Telegram. Sin esto no hay
        migración posible: Facebook solo ya es el 32,8\,\% del tráfico. Telegram ni
        siquiera figura en la hoja de ruta del producto y hay que pedirlo
        explícitamente.

  \item \textbf{El atajo \cod{/} en el chat.} Es el mecanismo declarado del CRM actual
        y lo usa el equipo todos los días, con decenas de atajos memorizados. Su
        ausencia no impide operar, pero degrada cada conversación.

  \item \textbf{Migración de contenido operativo.} Decenas de respuestas rápidas con
        variables que cambian de nombre, y decenas de etiquetas de negocio. Requiere un
        mapeo de variables, no una copia literal.

  \item \textbf{Mensajes programados.} Confirmado como necesario. Hoy existe en el CRM
        actual y no tiene equivalente en el nuevo.

  \item \textbf{Horario de atención por departamento} y \textbf{filtro por conexión}.

  \item \textbf{Autenticación en dos factores por usuario} y registro de
        \emph{última vez visto}.
\end{enumerate}

\subsection{Decisiones pendientes antes de estimar}

Tres secciones del CRM actual no tienen todavía una definición de si deben existir en
el nuevo. Mientras no se cierren, cualquier estimación de esfuerzo es incompleta:

\begin{center}
\begin{tabularx}{\textwidth}{@{}lL@{}}
\toprule
\textbf{Sección} & \textbf{Qué hay que definir} \\
\midrule
Campañas & Si el negocio sigue usando envíos masivos desde el CRM o los resuelve por
           otra vía \\
Informes & Si el \cod{/dashboard} del CRM nuevo reemplaza la reportería actual o hace
           falta construirla aparte \\
Tokens   & Qué integraciones externas dependen hoy de esas credenciales \\
\bottomrule
\end{tabularx}
\end{center}

Quedan descartadas por innecesarias: \textbf{Wäbot IA}, \textbf{Documentación} y
\textbf{Soporte}.

% =====================================================================
\section{Lo que no se pudo verificar}
\label{sec:limites}

Se declara explícitamente para que nadie tome por validado lo que no lo está.

\begin{itemize}
  \item \textbf{Arrastrar y soltar imágenes en el chat.} Arrastrar un archivo desde el
        escritorio es una acción del sistema operativo, fuera del alcance de la
        automatización usada. Debe probarse manualmente en el CRM nuevo: abrir un chat
        y arrastrar una imagen sobre la conversación.

  \item \textbf{El comportamiento de las respuestas rápidas con un departamento
        asignado.} El chat de prueba disponible en staging estaba ``Sin departamento'',
        por lo que solo se pudo comprobar el mensaje de ausencia. Falta verificar si,
        con departamento asignado, la lista aparece y si ofrece búsqueda y orden.

  \item \textbf{Las pantallas del CRM nuevo que no se recorrieron:} \cod{/dashboard},
        Chats internos, Mensajes destacados, Tickets bloqueados y las cinco de Control de
        acceso. Su existencia está confirmada; su funcionamiento no.

  \item \textbf{Si el CSAT del CRM nuevo se comporta igual al encenderlo.} Se verificó
        que la configuración existe y qué campos ofrece, pero está apagada: no se envió
        ninguna encuesta. No se sabe si repite por contacto, si permite dirigirla a un
        grupo ni si reabre la conversación al responderla.

  \item \textbf{La sección \emph{Resolución automática}} de la configuración del CRM
        nuevo, que aparece junto al CSAT y no se desplegó.

  \item \textbf{Las secciones del CRM actual no recorridas:} Campañas, Informes,
        Mensajes programados, Wäbot IA y Tokens. Solo se confirmó que existen en el
        menú.

  \item \textbf{Si los 93 permisos del perfil Admin incluyen alguno relacionado con
        visibilidad de contactos.} Se revisaron las diez categorías y sus contadores,
        pero no se desplegó cada permiso individual.
\end{itemize}

\begin{nota}
\textbf{Observación operativa.} Durante todo el recorrido, el CRM actual mostró de forma
recurrente un diálogo de \emph{autenticación en dos factores} que se superpone al
contenido. El equipo lo identifica como un defecto de presentación; se cierra pulsando
fuera del diálogo. No se interactuó con él en ningún momento.
\end{nota}

\end{document}
